\section{Reading a Leaderboard Number}
\label{app:leaderboard}

Section~\ref{sec:prior} quoted the two coupling costs the sources document
themselves: a tuned artifact is coupled to the serving model, and a Guide
tuned against a benchmark has to be measured on questions it has not seen.
The two systems atop DABStep's validated leaderboard make the point concrete
and disclose it. The first reports 89.95\% hard with a lightweight inference
model driving a helper library distilled offline by a heavyweight one, built
by ``validating them against ground truth answers'' and by testing candidate
functions ``against the ground truth of multiple interconnected tasks'',
with pitfalls extracted ``from the test data'' and injected into the
inference prompt~\cite{nvidia-kgmon}. The second, at 87.57\%, publishes under
the title \emph{Beyond Fine-tuning: Solving DABstep's Hard Mode with
Versioned Assets}~\cite{oceanbase-datapilot}; we could not retrieve that post
in full and make no claim about its mechanism. Nothing in the benchmark
forbids using its labelled dev split this way, and we allege no violation;
where the first system says ``the test data'' we read its own loose usage
rather than an admission. Such numbers answer a different question from
ours: how much of a labelled, enumerable question space can be compiled away
in advance.

Our gold reconstruction supplies one datum on the same point. Across the
2{,}199 published DABStep submissions, organisations submitting once achieve
a median hard-split accuracy of 19.6\%, while those submitting 21 or more
times achieve 54.2\%. That gradient is equally consistent with genuine
iterative improvement. But on a small-template benchmark with a public
leaderboard, a single accuracy carries little information about
generalisation unless the submission count and the artifact's provenance
are stated alongside it. Ours are: one sweep per model, no tuning against
the task set, a digest-pinned artifact frozen 35 commits before the first
run, and two leaderboard submissions, one per arm of the main contrast on
one of the four models, both made after the runs were complete and neither
followed by a change to the artifact. They exist to move that contrast onto
official grading and to audit our own golds (\appref{app:golds}), not
to place.

The comparable artifacts are unpublished: MotherDuck's and NVIDIA's are
fused to the test set and coupled to the model they were tuned against, and
OceanBase's is described only in a post we could not retrieve. Ours is in
the repository with a digest because a contract written from a vendor
manual has neither coupling. That
is the sense in which a data contract is a different object from a semantic
layer, and it costs us the accuracy gap rather than excusing it.

Two claims that might look distinctive here belong to others. Paired testing
across several models is not distinctive~\cite{semantic-paired,dbt-benchmark},
and paired McNemar is standard practice for comparing learners on one
run~\cite{dietterich}. The contribution is the decomposition.
